\section{Contribución esperada}
\label{sec:contribucion_esperada}

Para resolver los problemas encontrados, este proyecto propone el SISTEMA WEB DE INSCRIPCIÓN PARA EL MÓDULO TECNOLÓGICO PRODUCTIVO SAN JULIÁN BTH --- VERSIÓN 2.0. Es una página web sencilla que funciona bien tanto en celular como en computadora y que cuida que los datos se guarden de forma segura, evitando que dos personas ocupen el mismo cupo aunque entren al mismo tiempo. Es una mejora de la versión de 2025, porque agrega tres partes que faltaban: la recuperación de la contraseña sin necesidad de ir al colegio, el registro y revisión de los papeles entregados y el reparto automático de los estudiantes por colegio en cursos y paralelos, con un alcance importante para la gestión del BTH en todo el distrito, como señalan \textcite{sujovich2024}.

La página integra un formulario sencillo para inscribirse cuidando que los cupos no se dupliquen; una opción para recuperar la contraseña en pocos segundos con el número de carnet y la fecha de nacimiento, sin tener que hacer fila; un espacio para que secretaría registre y revise las libretas de 1ro a 3ro, las cédulas de los tutores y los certificados en el año que corresponde; y un reparto automático que ubica a cada estudiante en su curso y paralelo según su colegio, entregando al momento un comprobante para descargar con un código para verificar su validez.

Con esta mejora se espera un cambio muy positivo para el municipio. En lo social, se terminarían las filas de noche y se daría un acceso más justo y tranquilo a la educación técnica, al pasar de uno a tres días de espera a solo tres a cinco minutos por internet, y al permitir recuperar la contraseña en menos de dos minutos. En lo económico, se ahorrarían más de 20000 bolivianos al año entre pasajes y fotocopias para las familias, con un registro completo de todos los papeles, y se ampliaría la atención de 8 a 24 horas. Así, el Módulo se consolidaría como un ejemplo de modernización y buen manejo de la educación técnica en Santa Cruz.

Los detalles de cada una de estas dificultades y su diagnóstico entre 2022 y 2026 se presentan de forma ordenada en el Capítulo 2, donde se explica la realidad del proceso y se plantea el problema de investigación.
